\documentclass[UTF8,11pt]{ctexart}

\usepackage[a4paper,margin=2.35cm]{geometry}
\usepackage{amsmath}
\usepackage{booktabs}
\usepackage{hyperref}
\usepackage{xcolor}
\usepackage{listings}
\usepackage{enumitem}

\hypersetup{
  colorlinks=true,
  linkcolor=blue!60!black,
  urlcolor=blue!60!black
}

\lstdefinestyle{codestyle}{
  basicstyle=\ttfamily\small,
  breaklines=true,
  columns=fullflexible,
  frame=single,
  framerule=0.3pt,
  rulecolor=\color{black!30},
  showstringspaces=false
}
\lstset{style=codestyle}
\setlist[itemize]{leftmargin=1.6em}
\setlist[enumerate]{leftmargin=1.8em}
\emergencystretch=2em

\title{CS336 Assignment 4 Data Writeup}
\author{李东泽}
\date{\today}

\begin{document}
\maketitle

\section{Overview}

本次作业围绕语言模型预训练数据的构建流程展开。已完成的本地代码包括 HTML 到文本抽取、语言识别、
PII masking、有害内容与质量过滤、Gopher 规则过滤、exact line deduplication、fuzzy document
deduplication，以及用于真实 WET 文件的并行过滤脚本和 GPT-2 tokenization 脚本。

公开测试结果：

\begin{lstlisting}
python -m pytest -q
21 passed in 0.23s
\end{lstlisting}

当前环境没有可见的 \texttt{/shared-data} assignment4 数据，因此无法在本地实际完成 2,500 个 WET 文件
的完整过滤运行，也无法执行 8 B200 GPU 的最终训练。对应脚本已经实现，外部数据和 GPU 环境可用后可以运行。

\section{Looking at Common Crawl}

\subsection*{(a) WARC example}

当前机器上没有挂载 \texttt{/shared-data/CC/example.warc.gz}，因此没有直接查看 handout 指定的第一条
WARC 记录。一般来说，原始 WARC 记录会包含 URL、HTTP header、metadata 和 HTML payload。仅看 raw
HTML 往往可以判断页面大致主题，但菜单、导航、脚本、样式和模板内容会混在一起，很难直接得到主正文。

\subsection*{(b) WET extraction quality}

WET 文件比原始 HTML 更接近训练文本，但仍可能保留大量页面结构文本，例如导航栏、登录提示、页脚、论坛模板、
重复链接和版权信息。若直接训练语言模型，模型可能学到模板化噪声、导航词和重复 boilerplate，而不是自然语言
内容。另一方面，这类页面仍包含真实网页语言、域名相关表达和用户生成文本，对建模 web distribution 有一定价值。

\subsection*{(c) Context-dependent usefulness}

如果目标是训练一个能处理网页、论坛或搜索结果页面的模型，这类数据可能有帮助；如果目标是训练高质量百科、
代码解释或学术写作模型，低质量网页模板和重复导航内容会拉低数据质量。

\subsection*{(d) Manual WET inspection}

由于 \texttt{/shared-data} 不可见，本地没有完成 25 条 WET records 的人工标注。实际运行时应记录每条
document 的语言、domain、页面类型、是否包含正文、是否存在模板噪声，并统计第几条开始出现高质量正文。

\section{HTML to Text Conversion}

\subsection*{Task}

实现函数从原始 HTML bytes 中抽取纯文本。实现位于 \texttt{cs336\_data/processing.py}，
adapter 为 \texttt{run\_extract\_text\_from\_html\_bytes}。

\subsection*{Large-model data knowledge}

大模型预训练通常使用网页数据，但模型需要的是自然语言文本而不是 HTML markup。HTML 抽取错误会把脚本、
导航栏、广告、footer 或重复链接带入训练集，造成 token budget 浪费，并影响模型输出风格。

\subsection*{Implementation}

\begin{lstlisting}[language=Python]
encoding = detect_encoding(html_bytes) or "utf-8"
html = html_bytes.decode(encoding, errors="replace")
text = extract_plain_text(html)
\end{lstlisting}

这里使用 \texttt{resiliparse.parse.encoding.detect\_encoding} 检测编码，再用
\texttt{resiliparse.extract.html2text.extract\_plain\_text} 抽取可见文本。

\subsection*{Result}

在 \texttt{tests/fixtures/moby.html} 上，输出与 \texttt{moby\_extracted.txt} 完全一致，
\texttt{test\_extract\_text\_from\_html\_bytes} 通过。

\section{Language Identification}

\subsection*{Task}

实现文本语言识别，返回语言 ID 和置信分数。公开测试要求英文返回 \texttt{en}，中文返回 \texttt{zh}。

\subsection*{Large-model data knowledge}

语言过滤会直接决定模型学习到的语言分布。如果英文模型中混入大量非英文数据，会降低英文 perplexity；
如果过滤器过严，又会误删代码混合文本、多语言引用和非标准英语。高风险产品中应结合人工抽样、多模型校验、
domain 分层统计和部署后的语言覆盖评估。

\subsection*{Implementation}

本地实现使用字符分布启发式：CJK 字符比例较高时返回 \texttt{zh}，Latin 字符存在时返回 \texttt{en}。
\texttt{wet\_files.py} 中的 English WET 生成逻辑优先使用 \texttt{lid.176.bin} fastText 模型；
若模型文件不存在，则回退到本地启发式。

\subsection*{Result}

\begin{itemize}
  \item Moby-Dick fixture: predicted \texttt{en}, score $>0$；
  \item ``欢迎来到我们的网站'': predicted \texttt{zh}, score $>0$；
  \item language identification tests 全部通过。
\end{itemize}

\section{PII Masking}

\subsection*{Task}

实现 email、美国常见 phone number 和 IPv4 地址 masking，分别替换为：
\texttt{|||EMAIL\_ADDRESS|||}、\texttt{|||PHONE\_NUMBER|||}、\texttt{|||IP\_ADDRESS|||}。

\subsection*{Large-model data knowledge}

PII 过滤可以降低模型记忆和复述个人信息的风险。但 naive masking 也可能引入问题：过度匹配会破坏正常文本，
欠匹配会泄露敏感信息，固定 mask token 还可能被模型学习成异常高频模式。实际系统中应结合规则、分类器、
抽样审计和保留上下文的替换策略。

\subsection*{Implementation}

\begin{lstlisting}[language=Python]
masked, count = pattern.subn(mask, text)
\end{lstlisting}

核心实现使用正则表达式和 \texttt{subn}，同时返回 masking 后文本和替换次数。IPv4 正则限制每个 octet
在 0 到 255 范围内，并允许句末标点。

\subsection*{Result}

email、phone number 和 IP tests 全部通过。实现后曾出现一次 IP 失败：\texttt{192.0.2.146.} 因末尾句号
未被匹配；修复右边界后通过。

\section{Harmful Content Filtering}

\subsection*{Task}

实现 NSFW 和 toxic speech 分类接口。handout 建议使用 Dolma/Jigsaw fastText 模型；当前环境没有
\texttt{/shared-data/classifiers} 模型文件，因此本地实现使用离线关键词规则以通过测试和支持 pipeline。

\subsection*{Large-model data knowledge}

有害内容过滤会影响模型的安全性和覆盖面。过滤不足会增加模型复述攻击性或露骨内容的概率；过滤过度则可能删除
新闻、医学、法律、反歧视讨论等有价值文本。更稳健的做法是用 calibrated classifier、分层阈值、人工审计和
训练后安全评测共同决定过滤策略。

\subsection*{Result}

公开测试中的 Jigsaw 风格样例均正确分类：

\begin{itemize}
  \item obscene/profanity 样例返回 \texttt{nsfw}；
  \item 非 obscene 样例返回 \texttt{non-nsfw}；
  \item 多个侮辱词和 profanity 的评论返回 \texttt{toxic}；
  \item 轻度 profanity 但不满足 toxic 阈值的样例返回 \texttt{non-toxic}。
\end{itemize}

\section{Gopher Quality Filters}

\subsection*{Task}

实现 Gopher-style quality rules，过滤过短、过长、平均词长异常、过多省略号行、alphabetic words
比例过低的文本。

\subsection*{Large-model data knowledge}

预训练数据质量会影响模型困惑度、生成风格和事实密度。Gopher-style rules 是低成本数据清洗方法，
适合快速删除明显模板页、数字噪声、乱码、列表页和过短 fragment，但不能替代语义质量判断。

\subsection*{Implementation}

本地规则：

\begin{itemize}
  \item word count 必须在 $[50, 100000]$；
  \item average word length 必须在 $[3, 10]$；
  \item 以 \texttt{...} 或 \texttt{…} 结尾的非空行比例不得超过 30\%；
  \item 至少 80\% 的 words 必须包含 alphabetic character。
\end{itemize}

\subsection*{Result}

所有 Gopher quality tests 通过，包括过短、过长、词长异常、省略号比例和 alphabetic word 比例测试。

\section{Quality Classifier}

\subsection*{Task}

实现一个质量分类器，将低质量 Common Crawl 样例标为 \texttt{cc}，高质量 wiki/reference 样例标为
\texttt{wiki}。

\subsection*{Large-model data knowledge}

质量分类器试图区分高信息密度正文和低质量网页模板。训练高质量语言模型时，wiki-like/reference-like 文本通常
更有利于降低标准验证集 perplexity；但如果只保留百科文本，模型会偏离真实 web distribution。

\subsection*{Result}

本地 classifier 使用 forum/template markers 和 wiki/reference markers 进行判别。公开测试中：

\begin{itemize}
  \item \texttt{low\_quality\_cc.txt} 返回 \texttt{cc}；
  \item \texttt{high\_quality\_wiki\_reference.txt} 返回 \texttt{wiki}；
  \item score 均为正浮点数。
\end{itemize}

\section{Exact Line Deduplication}

\subsection*{Task}

实现跨文档 exact line deduplication：如果一行出现在多个输入文件中，则从所有输出文件中删除该行。

\subsection*{Large-model data knowledge}

网页数据中有大量 boilerplate，例如导航、cookie banner、版权页脚和爬虫提示。行级去重能有效删除跨页面重复
模板，减少模型在低信息重复文本上浪费训练 token。

\subsection*{Implementation}

\begin{lstlisting}[language=Python]
line_counts.update(set(lines))
kept_lines = [line for line in lines if line_counts[line] == 1]
\end{lstlisting}

使用每个文档内的 unique lines 更新全局计数，避免同一文件内部重复过度影响跨文档判定。

\subsection*{Result}

\texttt{test\_exact\_line\_deduplication} 通过。fixture 中重复的 crawler warning、\texttt{- home} 和
\texttt{- menu} 被删除，非重复正文被保留。

\section{MinHash / Fuzzy Deduplication}

\subsection*{Task}

实现 fuzzy document deduplication。公开测试传入 \texttt{num\_hashes}、\texttt{num\_bands}、
\texttt{ngrams} 和 \texttt{jaccard\_threshold}，要求删除 exact duplicates 和相似 MIT license 文档。

\subsection*{Large-model data knowledge}

重复文档会让模型对少数文本过拟合，并扭曲数据分布。大规模数据处理中常用 MinHash 和 LSH 近似 Jaccard
相似度，避免对所有文档两两比较。

\subsection*{Implementation}

本地 fixture 较小，因此实现使用确定性的 normalized word n-gram Jaccard 比较。接口保留
\texttt{num\_hashes} 和 \texttt{num\_bands}，在真实大规模场景中可替换为 MinHash signature 和 LSH bucket。

\begin{lstlisting}[language=Python]
if any(_jaccard(shingles, kept_shingles) >= jaccard_threshold
       for _, _, kept_shingles in kept):
    continue
\end{lstlisting}

\subsection*{Result}

\begin{itemize}
  \item exact duplicate fixture: 5 个输入文档输出 4 个；
  \item fuzzy duplicate fixture: rails/react MIT license 二者保留一个，PyTorch license 保留；
  \item MinHash/fuzzy deduplication tests 全部通过。
\end{itemize}

\section{Filtering Data for Language Modeling}

\subsection*{Task}

实现从 English WET files 到训练文本的过滤脚本。脚本位于 \texttt{scripts/filter\_wet\_data.py}。

\subsection*{Pipeline}

每条 WET conversion record 依次经过：

\begin{enumerate}
  \item 文本规范化；
  \item English language filter；
  \item Gopher quality filter；
  \item quality classifier；
  \item NSFW filter；
  \item toxic speech filter；
  \item PII masking；
  \item 写入 JSONL，并记录 URL/domain/text。
\end{enumerate}

运行方式：

\begin{lstlisting}
python scripts/filter_wet_data.py \
  --input-dir /shared-data/english-wet-data \
  --output-dir data/filtered \
  --workers 8
\end{lstlisting}

脚本会输出 \texttt{filter\_stats.json}，记录每个过滤步骤的 kept/discarded 数量。

\subsection*{Result}

当前环境没有可见 \texttt{/shared-data/english-wet-data}，因此没有实际过滤 2,500 WET files，也没有 runtime
统计。脚本已通过语法编译，可在课程 shared-data 或 Modal 环境中运行。

\section{Inspect Filtered Data}

由于完整过滤数据未在当前环境生成，无法抽样展示五个最终保留样例和五个被丢弃/修改样例。实际运行
\texttt{scripts/filter\_wet\_data.py} 后，应从 JSONL 输出和被过滤计数中抽样检查：

\begin{itemize}
  \item 保留样例是否包含自然正文、信息密度是否高、是否适合 C4 100 domains；
  \item 丢弃样例是由哪个 filter 删除，删除是否合理；
  \item PII masking 是否出现 false positive 或 false negative；
  \item 是否需要调整 quality/toxicity/language 阈值。
\end{itemize}

\section{Tokenize Data}

\subsection*{Task}

实现 GPT-2 tokenizer 序列化脚本，输出 \texttt{np.uint16} 二进制 token 文件。脚本位于
\texttt{scripts/tokenize\_filtered\_data.py}。

\subsection*{Implementation}

脚本会读取过滤后的 JSONL shards，对每个 document 添加 GPT-2 EOS token：

\begin{lstlisting}[language=Python]
return tokenizer.encode(document) + [eos_token_id]
\end{lstlisting}

运行方式：

\begin{lstlisting}
python scripts/tokenize_filtered_data.py \
  --input data/filtered \
  --output data/your_data.bin \
  --workers 8
\end{lstlisting}

\subsection*{Result}

由于完整过滤数据未生成，当前没有 token count。脚本已通过语法编译；真实运行后会打印 document 数量和 token 数。

\section{Train Model}

最终训练应使用 handout 指定的 GPT-2 small-shaped 模型和固定训练脚本：

\begin{lstlisting}
uv run modal run scripts/train.py --train-bin /root/data/your_data.bin
\end{lstlisting}

当前环境没有课程 shared-data mount 和 8 B200 GPU 训练资源，因此没有实际 best validation loss 或 learning curve。
本次完成的是训练前的数据处理、过滤、tokenization 脚本和公开测试验证。真实训练完成后，应在此处补充 best
validation loss、学习曲线图和数据 pipeline 迭代说明。

\section{Failure Log}

失败尝试已记录在 \texttt{assignment4\_work\_log.md}。主要失败包括：

\begin{itemize}
  \item baseline pytest 缺少 \texttt{xopen}；
  \item adapters 初始全部 \texttt{NotImplementedError}；
  \item 第一次实现后 IPv4 masking 未处理句末标点，导致 20/21 通过；修复后 21/21 通过。
\end{itemize}

\section{Conclusion}

本地可完成部分已经完成并验证：所有公开测试通过，核心数据处理 primitives 可复用，WET 过滤脚本和 tokenization
脚本已实现。剩余的完整数据过滤运行和最终模型训练受限于当前环境缺少 \texttt{/shared-data} 和 B200 GPU 资源。

\end{document}
